\documentclass{article} %210 mm × 297 mm
\usepackage[utf8]{inputenc}
\usepackage[a4paper,left=1.5cm, right=1.5cm, top=2cm, bottom=2cm]{geometry}
\usepackage{graphicx}
%\usepackage{blindtext}
\usepackage{multirow}
\usepackage{mhchem}
\usepackage{url}
\usepackage{subfigure}
\usepackage{amsfonts,latexsym} % para tener disponibilidad de diversos simbolos
\usepackage{enumerate}
%\usepackage{booktabs}
\usepackage{float}
%\usepackage{threeparttable}
\usepackage{array,colortbl}
%\usepackage{ifpdf}
%\usepackage{rotating}
\usepackage{cite}
\usepackage{stfloats}
\usepackage{url}
\usepackage{listings}
\usepackage{fancyhdr}

\usepackage{color}
%\usepackage{hyperref}
%\usepackage{wrapfig}
\usepackage{array}
%\usepackage{multirow}
%\usepackage{adjustbox}
%\usepackage{nccmath}
\usepackage[dvipsnames]{xcolor}

 
\newcommand{\titulo}{Präsenzübungsblatt 8; Statistik}
\newcommand{\nombre}{Lil'Jana}
\newcommand{\FCE}{\textbf{SoSe 2025}}
\newcommand{\dat}{\textbf{Tut 03.06.2025}}
\newcommand{\gre}[1]{\textcolor{green}{#1}}
\newcommand{\purp}[1]{\textcolor{purple}{#1}}
\newcommand{\cyan}[1]{\textcolor{cyan}{#1}}
\newcommand{\yel}[1]{\textcolor{yellow}{#1}}
\newcommand{\blu}[1]{\textcolor{blue}{#1}}
\newcommand{\red}[1]{\textcolor{red}{#1}}
\newcommand{\p}[0]{\mathbb{P}}
 

\usepackage{fancyhdr}
\pagestyle{fancy}
\fancyhead{}
\fancyfoot{}
\fancyhead[RO]{\small\dat}
\fancyhead[LO]{\small\FCE}
\fancyfoot[RO,LE] {\thepage}

\setcounter{page}{1} %Número de la página, la editorial lo cambiará

\newcommand{\kl}[1]{\{ #1 \}}

\usepackage[onehalfspacing]{setspace}
\begin{document}


\begin{tabular}{p{12cm}}
{{\Large\bf\titulo}}\\
\vspace{0.2cm}\\
 {\large\nombre}\\
\end{tabular}
\vspace{0.2cm}\\
\section{Aufgabe 6.1}
Foto, richtig in Abgabe 
\section{Aufgabe 6.2}
geg. $(x,y)$ uniform verteilt auf Menge $A$. \textcolor{orange}{d.h. sie haben eine gemeinsame Dichte $f(x,y) = \frac{\mathbb{I}_A(x,y)}{|A|}$}
\[
A = \kl{(x,y) \in \mathbb{R}^2 \mid -1 \leq c \leq 0, \quad x \leq y \leq 0} \cup \kl{(x,y) \in \mathbb{R}^2 \mid 0\leq x \leq 1, \quad x \leq y \leq 1}
\]
\subsection*{a)} korrekt 
\subsection{b)} korrekt 
\subsection{c)}
Vorzeichenfehler in Fall $0 \leq  x \leq 1$: 
\begin{align*}
    \fxx(x) &=\gre{ \int_{-1}^0 f_X (a) da} +\purp{ \int_0^x f_X (a) da}\\
    & = \gre{ \int_{-1}^0 1-a da} + \purp{\int_0^x 1-a da} \\
    &= \gre{[a -\frac{a^2}{2}]_{\yel{-1}}^\cyan{0}}  +  \purp{[a - \frac{a^2}{2}]^{\blu{x}}_{\red{0}}} \\
    &= \cyan{0} + \yel{(-1 + \frac{1}{2})} + \blu{x - \frac{x^2}{2}} + \red{0} \\
    & = \yel{\frac{1}{2}} + \blu{x - \frac{x^2}{2}} \\
    & = x + \frac{3-x^2}{2}
\end{align*}
Flüchtigkeitsfehler in Fall $0 \leq y \leq 1$
\begin{align*}
    
\end{align*}
\subsection*{d)}
\subsubsection*{(ii)}
Verteilungsfunktion. Hier reicht es, wenn man sich nur einen Wert ausucht, weil das $X > 0$ direkt eine ganze Fläche ist. 
\begin{align*}
    \p(X > 0, Y \leq 0) & = 0\\ \\
    \p (X>0) &= 1 - \p(X \leq 0) \\
    &= 1 - \frac{1}{2} \\
    &= \frac{1}{2} \\ \\
    \p(Y \leq 0) &= \frac{1}{2}
\end{align*}
\section{Aufgabe 6.3}
Für Wartezeiten benutzt man eig. immer die Exponentialverteilung. 
\subsection*{a)}
Korrekt
\subsection{b)}
Wie nannt man diese Eigenschaft? Gedächtnislosigkeit. Die kommt auch bei der geomterischen Verteilung vor.\\
\[
\p(X \geq 40 \mid x \geq  20) = \frac{\p(x \geq 40, X \geq 20)}{\p(X\geq 20)} = \frac{\p(X \geq 40)}{\p(X \geq 20)} = \frac{e^{-4}}{e^{-2}} = e^{-2}
\]

\subsection*{c)}
$Y \sim Exp(\frac{1}{10})$ Fahrzeit, $\p(X+Y \leq 30)$, $\alpha = \frac{1}{10}$
\begin{align*}
     \p(X+Y \leq Z) &= \int_\substack{x+y \leq z \\ x \geq 0, y \geq 0} \underbrace{\alpha e^{-\alpha x } \alpha e^{-\alpha y}}_{gem. Dichte} d (x,y) \\
     &= \int_0^z  ( \int_0^{z-x} \alpha^2 e^{- \alpha x} e^{- \alpha y} dy) dx \\
     &= \alpha^2 \int_0^z e^{- \alpha x} [-\frac{1}{\alpha} a^{\alpha}y]^{z-x}_0 dx \\
     &= - \alpha \int_0^z e^{- \alpha x} (e^{- \alpha z} e^{\alpha x} -1) dx \\
     &= - \alpha \int_0^z e^{- \alpha z} - e^{- \alpha x} dx \\
     &= - \alpha [x e^{- \alpha z} + \frac{1}{\alpha} e^{- \alpha x}]^z_0 \\
     &= - \alpha (z e^{- \alpha z}+ \frac{1}{\alpha} e^{-\alpha z} + \frac{1}{\alpha}) \\
     &= 1 - e^{-\alpha z} (\alpha z +1)
\end{align*}
$ (\alpha z +1)$ oder $ (-\alpha z +1)$? Eher ersteres, nachrechnen

\section{Präsenz}
\subsection{Aufgabe}
Reminder: \\
Erwartungswert: \( \mathbb{E}[X] = \underbrace{\sum_{j \in \mathbb{Z}} j \p(X=j)}_{< \infty}  \) als ''gewichteter Mittelwert''. \\
Varianz \(Var(X) = \mathbb{E}[(X-\mathbb{E}[X])^2]\) \\
\textcolor{pink}{Frage: Müssen die $\p$s von $X$ zusammen 1 ergeben? Ja müssen sie.} \\
\textcolor{brown}{
Wert für $X=0$ berechnen: \\
\(\p(X=-1) + \p(X=1) + \p(X=0) + \p(X=2) + \p(X=-2) = 1\) also 
\[
\frac{p}{4} + \frac{p}{4} + \p(X=0) + \frac{1}{4} + \frac{1}{4} = 1
\]
\[
\frac{2p}{4} + \p(X=0) + \frac{2}{4}  = 1
\]
\[
\frac{2p + 2}{4} + \p(X=0)  = 1
\]
\[
 \p(X=0)  = 1 - \frac{2p + 2}{4}= \frac{4}{4} - \frac{2p + 2}{4} = \frac{4-2p - 2}{4} = \frac{2-2p}{4} = \frac{1-p}{2}
\]
}
\subsubsection*{a)}
Foto und so 

\end{document}